\Anexo{Repositorio de Código y Configuraciones Técnicas}

Este anexo detalla los enlaces a los repositorios donde se encuentra el código fuente del Asistente Virtual Inteligente, así como la configuración base de las variables de entorno para su correcto funcionamiento.

\textbf{Repositorio de Código Fuente}
\begin{itemize}
    \item \textbf{Asistente Virtual (Código Fuente):} \url{https://github.com/EmilioEJ/Asistente_Virtual}
    \item \textbf{Documento de Tesis (LaTeX):} \url{https://github.com/EmilioEJ/TFG_EmilioESPINOZA}
\end{itemize}

\textbf{Configuración Base (Archivo .env)}
A continuación se presenta un extracto de las variables de entorno necesarias para levantar el backend en FastAPI:

\begin{verbatim}
# Configuración del Servidor y CORS
PORT=8000
FRONTEND_URL=http://localhost:5173

# Modelos y APIs de IA
ELEVENLABS_API_KEY=tu_clave_aqui
ELEVENLABS_VOICE_ID=voz_valeria
GROQ_API_KEY=tu_clave_aqui
WHISPER_MODEL=whisper-large-v3

# Base de Datos Vectorial Local (ChromaDB)
CHROMA_DB_PATH=./chroma_db
EMBEDDING_MODEL=sentence-transformers/paraphrase-multilingual-MiniLM-L12-v2

# Modelo LLM Local / API
NEMOTRON_ENDPOINT=http://localhost:11434/api/generate
\end{verbatim}

\Anexo{Instrumento de Evaluación de Usabilidad}

Para la evaluación del sistema, se utilizó un cuestionario basado en Google Forms distribuido a 13 estudiantes de la carrera. A continuación, se presenta la estructura de las preguntas formuladas:

\begin{itemize}
    \item \textbf{Sección 1: Perfil Técnico del Evaluador}
    \begin{enumerate}
        \item ¿Cuál es tu nivel de avance actual dentro de la carrera de TI?
        \item En los últimos 6 meses, ¿con qué frecuencia has interactuado con modelos de IA generativa para actividades académicas? (Escala 1-5)
        \item ¿Cuál es tu nivel de conocimiento técnico en renderización gráfica 3D o desarrollo interactivo? (Escala 1-5)
    \end{enumerate}
    \item \textbf{Sección 2: Usabilidad de la Interfaz y Arquitectura}
    \begin{enumerate}
        \setcounter{enumi}{3}
        \item Considero que las distintas funciones del sistema están bien integradas.
        \item Encontré que el sistema es innecesariamente complejo de utilizar.
        \item La información académica que provee el asistente es clara, precisa y fácil de entender.
        \item El sistema muestra mensajes de error claros o me guía adecuadamente ante fallos.
    \end{enumerate}
    \item \textbf{Sección 3: Inteligencia Conversacional y NLP}
    \begin{enumerate}
        \setcounter{enumi}{7}
        \item Al inicio, el asistente explicó claramente su propósito y límites.
        \item El asistente virtual comprendió mis intenciones y consultas técnicas con precisión.
        \item Las respuestas de la IA fueron útiles y directamente relevantes.
        \item Cuando la IA no entendió algo, fue capaz de reconducir la conversación amigablemente.
    \end{enumerate}
    \item \textbf{Sección 4: Evaluación de la Interfaz del Avatar 3D}
    \begin{enumerate}
        \setcounter{enumi}{11}
        \item Antropomorfismo. (1: Artificial - 5: Natural)
        \item Animacidad. (1: Estático - 5: Vivo)
        \item Afinidad visual y empatía. (1: Desagradable - 5: Agradable)
        \item Inteligencia percibida. (1: Incompetente - 5: Competente)
    \end{enumerate}
    \item \textbf{Sección 5: Visión Artificial y Aceptación Tecnológica}
    \begin{enumerate}
        \setcounter{enumi}{15}
        \item Percibí una latencia aceptable entre la captura y la respuesta renderizada.
        \item El uso de la cámara me pareció poco invasivo y confío en el manejo de privacidad.
        \item Utilizar este asistente aumentaría mi eficiencia para encontrar información.
        \item Basado en esta prueba, preferiría utilizar este asistente en el futuro.
    \end{enumerate}
    \item \textbf{Sección 6: Diagnóstico Técnico}
    \begin{enumerate}
        \setcounter{enumi}{19}
        \item ¿Qué mejora técnica priorizarías antes de lanzar este asistente a toda la universidad?
    \end{enumerate}
\end{itemize}

\Anexo{Formulario de Consentimiento Informado}

El siguiente formulario fue presentado a todos los participantes antes de la evaluación, indicando explícitamente el uso de la cámara web sin almacenamiento, asegurando la privacidad, el anonimato de los datos y el carácter voluntario de la prueba.

\begin{figure}[htbp]
    \centering
    \includegraphics[width=0.95\textwidth]{Extras/Anexo_Consentimiento.png}
    \caption{Formulario de consentimiento informado.}
\end{figure}

\Anexo{Aprobación de Abstract}

A continuación se adjunta la evidencia de la aprobación del documento de resumen ejecutivo en su versión en inglés (Abstract).

\begin{figure}[htbp]
    \centering
    \includegraphics[width=0.85\textwidth]{Extras/Anexo_Abstract.png}
    \caption{Aprobación del Abstract.}
\end{figure}